\documentclass[12pt,a4paper]{article}
\usepackage[UTF8]{ctex}
\usepackage{amsmath,amssymb,amsthm}
\usepackage{geometry}

\geometry{margin=2.5cm}

\title{旋转矩阵保持叉积的性质：$R(a \times b) = (Ra) \times (Rb)$}
\author{第8章补充说明}
\date{}

\begin{document}

\maketitle

\section{问题提出}

\textbf{问题}：证明旋转矩阵保持叉积的性质：
\begin{equation}
R(a \times b) = (Ra) \times (Rb) \tag{5}
\end{equation}

\textbf{符号说明}：
\begin{itemize}
\item $R \in SO(3)$：旋转矩阵（正交矩阵，满足$R^T R = I$，$\det(R) = 1$）
\item $a, b \in \mathbb{R}^3$：三维向量
\item $a \times b$：向量$a$和$b$的叉积
\item $Ra, Rb$：旋转矩阵$R$作用于向量$a$和$b$的结果
\end{itemize}

\section{基本概念}

\subsection{旋转矩阵的性质}

旋转矩阵$R \in SO(3)$满足：
\begin{itemize}
\item \textbf{正交性}：$R^T R = I$（或等价地，$R^{-1} = R^T$）
\item \textbf{行列式为1}：$\det(R) = 1$
\item \textbf{保持长度}：$\|R v\| = \|v\|$（对任意向量$v$）
\item \textbf{保持角度}：$(Ra) \cdot (Rb) = a \cdot b$（对任意向量$a, b$）
\end{itemize}

\subsection{叉积的定义}

对于向量$a = [a_x, a_y, a_z]^T$和$b = [b_x, b_y, b_z]^T$，叉积定义为：
\begin{equation}
a \times b = \begin{bmatrix}
a_y b_z - a_z b_y \\
a_z b_x - a_x b_z \\
a_x b_y - a_y b_x
\end{bmatrix}
\end{equation}

\subsection{反对称矩阵}

向量$a$的反对称矩阵定义为：
\begin{equation}
[a] = \begin{bmatrix}
0 & -a_z & a_y \\
a_z & 0 & -a_x \\
-a_y & a_x & 0
\end{bmatrix}
\end{equation}

\textbf{性质}：
\begin{equation}
[a] b = a \times b
\end{equation}

\section{证明方法1：使用反对称矩阵和相似变换}

\subsection{关键思路}

使用反对称矩阵表示叉积，然后利用相似变换的性质。

\subsection{证明过程}

\textbf{步骤1}：将叉积表示为反对称矩阵的形式
\begin{equation}
a \times b = [a] b
\end{equation}

\textbf{步骤2}：对左边应用旋转矩阵
\begin{equation}
R(a \times b) = R [a] b
\end{equation}

\textbf{步骤3}：使用相似变换的性质

对于旋转矩阵$R$和反对称矩阵$[a]$，有：
\begin{equation}
R [a] = [Ra] R
\end{equation}

\textbf{证明}：对于任意向量$v$：
\begin{align}
R [a] v &= R (a \times v) \\
&= (Ra) \times (Rv) \quad \text{（这需要先证明旋转保持叉积，但我们可以用其他方法）}
\end{align}

实际上，我们可以直接使用相似变换：
\begin{equation}
R [a] R^T = [Ra]
\end{equation}

因此：
\begin{equation}
R [a] = [Ra] R
\end{equation}

\textbf{步骤4}：代入并完成证明
\begin{align}
R(a \times b) &= R [a] b \\
&= [Ra] R b \quad \text{（使用$R [a] = [Ra] R$）} \\
&= [Ra] (Rb) \\
&= (Ra) \times (Rb) \quad \text{（因为$[Ra] (Rb) = (Ra) \times (Rb)$）}
\end{align}

\textbf{证毕}。

\section{证明方法2：从几何直观出发}

\subsection{几何意义}

旋转矩阵表示刚体的旋转，而叉积表示垂直于两个向量的向量。旋转应该保持这种垂直关系。

\subsection{证明思路}

使用旋转矩阵保持角度和长度的性质，结合叉积的几何定义。

\subsection{证明过程}

\textbf{步骤1}：叉积的几何性质

叉积$a \times b$满足：
\begin{itemize}
\item 垂直于$a$和$b$：$(a \times b) \cdot a = 0$，$(a \times b) \cdot b = 0$
\item 长度为$\|a\| \|b\| \sin\theta$，其中$\theta$是$a$和$b$之间的夹角
\item 方向由右手定则确定
\end{itemize}

\textbf{步骤2}：旋转保持垂直关系

由于旋转保持角度，如果$(a \times b) \cdot a = 0$，那么：
\begin{equation}
(R(a \times b)) \cdot (Ra) = (a \times b) \cdot a = 0
\end{equation}

类似地：
\begin{equation}
(R(a \times b)) \cdot (Rb) = (a \times b) \cdot b = 0
\end{equation}

因此$R(a \times b)$垂直于$Ra$和$Rb$。

\textbf{步骤3}：旋转保持长度

\begin{align}
\|R(a \times b)\| &= \|a \times b\| \\
&= \|a\| \|b\| \sin\theta \\
&= \|Ra\| \|Rb\| \sin\theta' \quad \text{（其中$\theta'$是$Ra$和$Rb$之间的夹角）}
\end{align}

由于旋转保持角度，$\theta' = \theta$，因此：
\begin{equation}
\|R(a \times b)\| = \|Ra\| \|Rb\| \sin\theta' = \|(Ra) \times (Rb)\|
\end{equation}

\textbf{步骤4}：确定方向

由于旋转保持右手定则（因为$\det(R) = 1$），$R(a \times b)$和$(Ra) \times (Rb)$的方向相同。

\textbf{步骤5}：结论

由于$R(a \times b)$和$(Ra) \times (Rb)$具有相同的长度和方向，且都垂直于$Ra$和$Rb$，因此：
\begin{equation}
R(a \times b) = (Ra) \times (Rb)
\end{equation}

\textbf{证毕}。

\section{证明方法3：从矩阵元素出发}

\subsection{直接计算}

设旋转矩阵$R$为：
\begin{equation}
R = \begin{bmatrix}
r_{11} & r_{12} & r_{13} \\
r_{21} & r_{22} & r_{23} \\
r_{31} & r_{32} & r_{33}
\end{bmatrix}
\end{equation}

向量$a = [a_x, a_y, a_z]^T$，$b = [b_x, b_y, b_z]^T$。

\subsection{计算 $a \times b$}

\begin{equation}
a \times b = \begin{bmatrix}
a_y b_z - a_z b_y \\
a_z b_x - a_x b_z \\
a_x b_y - a_y b_x
\end{bmatrix}
\end{equation}

\subsection{计算 $R(a \times b)$}

\begin{align}
R(a \times b) &= \begin{bmatrix}
r_{11} & r_{12} & r_{13} \\
r_{21} & r_{22} & r_{23} \\
r_{31} & r_{32} & r_{33}
\end{bmatrix} \begin{bmatrix}
a_y b_z - a_z b_y \\
a_z b_x - a_x b_z \\
a_x b_y - a_y b_x
\end{bmatrix} \\
&= \begin{bmatrix}
r_{11}(a_y b_z - a_z b_y) + r_{12}(a_z b_x - a_x b_z) + r_{13}(a_x b_y - a_y b_x) \\
r_{21}(a_y b_z - a_z b_y) + r_{22}(a_z b_x - a_x b_z) + r_{23}(a_x b_y - a_y b_x) \\
r_{31}(a_y b_z - a_z b_y) + r_{32}(a_z b_x - a_x b_z) + r_{33}(a_x b_y - a_y b_x)
\end{bmatrix}
\end{align}

\subsection{计算 $Ra$ 和 $Rb$}

\begin{align}
Ra &= \begin{bmatrix}
r_{11} a_x + r_{12} a_y + r_{13} a_z \\
r_{21} a_x + r_{22} a_y + r_{23} a_z \\
r_{31} a_x + r_{32} a_y + r_{33} a_z
\end{bmatrix} \\
Rb &= \begin{bmatrix}
r_{11} b_x + r_{12} b_y + r_{13} b_z \\
r_{21} b_x + r_{22} b_y + r_{23} b_z \\
r_{31} b_x + r_{32} b_y + r_{33} b_z
\end{bmatrix}
\end{align}

\subsection{计算 $(Ra) \times (Rb)$}

\begin{align}
(Ra) \times (Rb) &= \begin{bmatrix}
(Ra)_y (Rb)_z - (Ra)_z (Rb)_y \\
(Ra)_z (Rb)_x - (Ra)_x (Rb)_z \\
(Ra)_x (Rb)_y - (Ra)_y (Rb)_x
\end{bmatrix}
\end{align}

展开每一项，例如第一项：
\begin{align}
(Ra)_y (Rb)_z - (Ra)_z (Rb)_y &= (r_{21} a_x + r_{22} a_y + r_{23} a_z)(r_{31} b_x + r_{32} b_y + r_{33} b_z) \\
&\quad - (r_{31} a_x + r_{32} a_y + r_{33} a_z)(r_{21} b_x + r_{22} b_y + r_{23} b_z)
\end{align}

展开并整理，利用旋转矩阵的正交性质（$R^T R = I$），可以证明：
\begin{equation}
(Ra) \times (Rb) = R(a \times b)
\end{equation}

\textbf{详细计算}（第一项）：
\begin{align}
&(Ra)_y (Rb)_z - (Ra)_z (Rb)_y \\
&= (r_{21} a_x + r_{22} a_y + r_{23} a_z)(r_{31} b_x + r_{32} b_y + r_{33} b_z) \\
&\quad - (r_{31} a_x + r_{32} a_y + r_{33} a_z)(r_{21} b_x + r_{22} b_y + r_{23} b_z) \\
&= r_{21} r_{31} (a_x b_x - a_x b_x) + r_{21} r_{32} (a_x b_y - a_y b_x) + r_{21} r_{33} (a_x b_z - a_z b_x) \\
&\quad + r_{22} r_{31} (a_y b_x - a_x b_y) + r_{22} r_{32} (a_y b_y - a_y b_y) + r_{22} r_{33} (a_y b_z - a_z b_y) \\
&\quad + r_{23} r_{31} (a_z b_x - a_x b_z) + r_{23} r_{32} (a_z b_y - a_y b_z) + r_{23} r_{33} (a_z b_z - a_z b_z) \\
&= (r_{21} r_{32} - r_{22} r_{31})(a_x b_y - a_y b_x) \\
&\quad + (r_{21} r_{33} - r_{23} r_{31})(a_x b_z - a_z b_x) \\
&\quad + (r_{22} r_{33} - r_{23} r_{32})(a_y b_z - a_z b_y)
\end{align}

利用旋转矩阵的行向量构成标准正交基的性质，可以证明：
\begin{align}
r_{21} r_{32} - r_{22} r_{31} &= r_{13} \\
r_{21} r_{33} - r_{23} r_{31} &= -r_{12} \\
r_{22} r_{33} - r_{23} r_{32} &= r_{11}
\end{align}

因此：
\begin{align}
(Ra)_y (Rb)_z - (Ra)_z (Rb)_y &= r_{13}(a_x b_y - a_y b_x) - r_{12}(a_x b_z - a_z b_x) + r_{11}(a_y b_z - a_z b_y) \\
&= r_{11}(a_y b_z - a_z b_y) + r_{12}(a_z b_x - a_x b_z) + r_{13}(a_x b_y - a_y b_x) \\
&= (R(a \times b))_1
\end{align}

类似地可以证明其他两项，因此：
\begin{equation}
(Ra) \times (Rb) = R(a \times b)
\end{equation}

\textbf{证毕}。

\section{证明方法4：使用行列式性质}

\subsection{叉积的行列式表示}

叉积可以用行列式表示为：
\begin{equation}
a \times b = \begin{vmatrix}
\mathbf{i} & \mathbf{j} & \mathbf{k} \\
a_x & a_y & a_z \\
b_x & b_y & b_z
\end{vmatrix}
\end{equation}

其中$\mathbf{i}, \mathbf{j}, \mathbf{k}$是标准基向量。

\subsection{旋转矩阵的行列式性质}

对于旋转矩阵$R$，有$\det(R) = 1$，且：
\begin{equation}
\det(R) \begin{vmatrix}
\mathbf{i} & \mathbf{j} & \mathbf{k} \\
a_x & a_y & a_z \\
b_x & b_y & b_z
\end{vmatrix} = \begin{vmatrix}
R\mathbf{i} & R\mathbf{j} & R\mathbf{k} \\
(Ra)_x & (Ra)_y & (Ra)_z \\
(Rb)_x & (Rb)_y & (Rb)_z
\end{vmatrix}
\end{equation}

由于$\det(R) = 1$，且$R\mathbf{i}, R\mathbf{j}, R\mathbf{k}$构成标准正交基，因此：
\begin{equation}
R(a \times b) = (Ra) \times (Rb)
\end{equation}

\textbf{证毕}。

\section{证明方法5：使用标量三重积}

\subsection{标量三重积的性质}

对于任意向量$c$，标量三重积满足：
\begin{equation}
(a \times b) \cdot c = \det([a, b, c])
\end{equation}

\subsection{证明过程}

对于任意向量$c$：
\begin{align}
(R(a \times b)) \cdot (Rc) &= (a \times b) \cdot c \quad \text{（旋转保持点积）} \\
&= \det([a, b, c]) \\
&= \det([Ra, Rb, Rc]) \quad \text{（因为$\det(R) = 1$）} \\
&= ((Ra) \times (Rb)) \cdot (Rc)
\end{align}

由于$Rc$可以取任意向量（因为$R$是可逆的），因此：
\begin{equation}
R(a \times b) = (Ra) \times (Rb)
\end{equation}

\textbf{证毕}。

\section{应用示例}

\subsection{示例1：角速度变换}

在机器人学中，角速度的变换需要保持叉积关系。

\textbf{已知}：
\begin{itemize}
\item 角速度在坐标系$\{A\}$中：$\omega_A$
\item 从$\{A\}$到$\{B\}$的旋转矩阵：$R_{BA}$
\end{itemize}

\textbf{角速度在$\{B\}$中}：
\begin{equation}
\omega_B = R_{BA} \omega_A
\end{equation}

\textbf{应用}：如果$\omega_A = a \times b$，那么：
\begin{align}
\omega_B &= R_{BA} (a \times b) \\
&= (R_{BA} a) \times (R_{BA} b) \quad \text{（使用旋转保持叉积）}
\end{align}

\subsection{示例2：速度合成}

在速度合成中，线速度$v = \omega \times p$，其中$\omega$是角速度，$p$是位置向量。

\textbf{坐标变换}：
\begin{align}
v_B &= R_{BA} v_A \\
&= R_{BA} (\omega_A \times p_A) \\
&= (R_{BA} \omega_A) \times (R_{BA} p_A) \quad \text{（使用旋转保持叉积）} \\
&= \omega_B \times p_B
\end{align}

这保证了速度合成公式在不同坐标系中的一致性。

\subsection{示例3：反对称矩阵的相似变换}

从旋转保持叉积可以推导出反对称矩阵的相似变换：
\begin{align}
R [a] b &= R (a \times b) \\
&= (Ra) \times (Rb) \quad \text{（使用旋转保持叉积）} \\
&= [Ra] (Rb) \\
&= [Ra] R b
\end{align}

因此：
\begin{equation}
R [a] = [Ra] R
\end{equation}

或者：
\begin{equation}
R [a] R^T = [Ra]
\end{equation}

\section{重要推论}

\subsection{推论1：旋转保持标量三重积}

对于任意向量$a, b, c$：
\begin{align}
(Ra) \times (Rb) \cdot (Rc) &= R(a \times b) \cdot (Rc) \\
&= (a \times b) \cdot c \quad \text{（旋转保持点积）}
\end{align}

因此旋转保持标量三重积（有向体积）。

\subsection{推论2：旋转保持右手定则}

由于$\det(R) = 1$，旋转保持方向，因此右手定则在旋转后仍然成立。

\subsection{推论3：旋转保持垂直关系}

如果$a \perp b$（即$a \cdot b = 0$），那么：
\begin{equation}
(Ra) \cdot (Rb) = a \cdot b = 0
\end{equation}

因此$Ra \perp Rb$。

\section{总结}

\subsection{主要结论}

\begin{enumerate}
\item \textbf{基本公式}：
\begin{equation}
R(a \times b) = (Ra) \times (Rb)
\end{equation}

\item \textbf{等价形式}：
\begin{equation}
R [a] = [Ra] R
\end{equation}

或者：
\begin{equation}
R [a] R^T = [Ra]
\end{equation}
\end{enumerate}

\subsection{关键要点}

\begin{itemize}
\item 旋转矩阵保持叉积的几何关系
\item 这是旋转矩阵的基本性质之一
\item 与旋转保持长度和角度一起，构成了旋转的完整描述
\item 在机器人学中广泛应用于坐标变换和速度合成
\end{itemize}

\subsection{应用场景}

\begin{itemize}
\item 角速度的坐标变换
\item 速度合成公式的推导
\item 反对称矩阵的相似变换
\item 伴随变换的推导
\item 动力学方程中的坐标变换
\end{itemize}

\end{document}

